\documentclass[11pt]{article}
\usepackage[margin=1in]{geometry}
\usepackage{amsmath,amssymb,amsthm}
\usepackage{graphicx}
\usepackage[hidelinks]{hyperref}

\newtheorem{theorem}{Theorem}

\newcommand{\R}{\mathbb{R}}
\newcommand{\rhoinf}{\rho}
\newcommand{\rhobar}{\overline{\rho}}

\title{A Bounded Collapsed Almost Uncollapsed Map}
\author{Automated discovery packet}
\date{June 14, 2026}

\begin{document}
\maketitle

\section*{Source question}

Braga's paper \cite{Braga} asks Problem 2.6:
``Is there a map \(X\to Y\) which is collapsed, almost uncollapsed and
bounded (in particular not solvent), for some Banach spaces \(X\) and \(Y\)?''

\begin{center}
\includegraphics[width=\textwidth]{figures/open_problem_crop.png}
\end{center}

\section*{Definitions used}

For a map \(f:(M,d)\to(N,\partial)\), the source defines
\[
  \rhoinf_f(t)=\inf\{\partial(f(x),f(y)):\ d(x,y)\geq t\}
\]
and
\[
  \rhobar_f(t)=\inf\{\partial(f(x),f(y)):\ d(x,y)=t\}.
\]
The map is uncollapsed if \(\rhoinf_f(t)>0\) for some \(t>0\), and almost
uncollapsed if \(\rhobar_f(t)>0\) for some \(t>0\). In the terminology of the
paper, collapsed means not uncollapsed.

\section*{Proof Intuition}

The exact-distance requirement in almost uncollapsedness is much weaker than
the large-distance requirement in uncollapsedness. A periodic map into the
unit circle separates all pairs at one chosen distance, for instance distance
\(\pi\), but repeats values at arbitrarily large distances. This gives
almost uncollapsedness and collapsedness at the same time. Since the circle is
bounded, the same map is automatically not solvent. If one wants the domain
to be infinite-dimensional, the same idea can be made set-theoretic: identify
points modulo a large-period cyclic subgroup and send different cosets to
different unit vectors.

\section*{Result}

\begin{theorem}
There are Banach spaces \(X,Y\) and a bounded map \(f:X\to Y\) which is
collapsed and almost uncollapsed. In particular, Problem 2.6 of \cite{Braga}
has an affirmative answer. Moreover, if the problem is read as asking for an
infinite-dimensional domain, \(X\) can be chosen infinite-dimensional.
\end{theorem}

\begin{proof}
Take \(X=\R\), take \(Y=\R^2\) with its Euclidean norm, and define
\[
       f(x)=(\cos x,\sin x).
\]
The map is bounded, since \(\|f(x)\|_2=1\) for every \(x\in\R\).

It is almost uncollapsed. Indeed, if \(|x-y|=\pi\), then \(y=x+\pi\) or
\(y=x-\pi\), and hence \(f(y)=-f(x)\). Therefore
\[
       \|f(x)-f(y)\|_2=2
\]
for every pair \(x,y\in\R\) with \(|x-y|=\pi\). Thus
\[
       \rhobar_f(\pi)=2>0.
\]

It is collapsed. Let \(t>0\). Choose an integer \(k\geq1\) such that
\(2\pi k\geq t\). Then \(|2\pi k-0|\geq t\), while
\[
       f(2\pi k)=f(0).
\]
Consequently \(\rhoinf_f(t)=0\) for every \(t>0\). Thus \(f\) is not
uncollapsed, i.e. it is collapsed in the source terminology.

Finally, \(f\) is not solvent. Its image has diameter at most \(2\), so for
any \(n>2\) no interval of domain distances can force all image distances to
be greater than \(n\). This proves the claimed bounded collapsed almost
uncollapsed example.

For the infinite-dimensional variant, let \(X\) be any infinite-dimensional
Banach space and choose \(u\in X\) with \(\|u\|=1\). Let
\[
       A=2\pi\mathbb{Z}u
\]
and let \(I=X/A\) be the set of cosets of the additive subgroup \(A\). Put
\(Y=\ell_2(I)\), with its standard unit vectors \((e_i)_{i\in I}\), and define
\[
       F(x)=e_{x+A}.
\]
This map is bounded, since every value has norm \(1\). If \(\|x-y\|=\pi\) and
\(x+A=y+A\), then \(x-y=2\pi k u\) for some \(k\in\mathbb{Z}\), so
\(\|x-y\|=2\pi |k|\), impossible. Hence exact-distance-\(\pi\) pairs always
belong to different cosets, and therefore
\[
       \|F(x)-F(y)\|_2=\sqrt{2}
       \qquad(\|x-y\|=\pi).
\]
Thus \(\rhobar_F(\pi)=\sqrt{2}>0\). On the other hand, for every \(t>0\) one
can choose \(k\geq1\) with \(2\pi k\geq t\), and then
\[
       \|2\pi k u-0\|\geq t,\qquad F(2\pi k u)=F(0).
\]
So \(\rhoinf_F(t)=0\) for every \(t>0\), and \(F\) is collapsed. Its image
has diameter \(\sqrt{2}\), so it is not solvent. This variant is not claimed
to be continuous; it is included only to remove any possible finite-dimensional
reading concern.
\end{proof}

\section*{Verification notes}

The first proof uses only the source definitions and the elementary identities
\(f(x+\pi)=-f(x)\) and \(f(x+2\pi k)=f(x)\). The infinite-dimensional variant
uses only cosets of the cyclic subgroup \(2\pi\mathbb{Z}u\). The script
\texttt{code/check\_circle\_map.py} performs a finite numerical smoke check
of these identities and the boundedness normalization; it is not needed for
the proof. The script \texttt{code/make\_open\_problem\_crop.py} recreates
the source screenshot from \texttt{source\_paper.pdf}.

Verifier focus: confirm that the paper uses ``collapsed'' as the negation of
``uncollapsed'', as indicated in the surrounding text before Proposition 2.5,
and that Problem 2.6 does not impose any additional regularity such as being a
cocycle or a coarse map. The finite-dimensional displayed example is in fact
Lipschitz, but it is not coarse-embedding-like because it is bounded. The
infinite-dimensional variant answers only the unrestricted-map reading.

\section*{Novelty and literature check}

Local duplicate check on June 14, 2026 searched
\texttt{registry\_index.tsv}, \texttt{solutions/index.tsv},
\texttt{attempts/index.tsv}, and \texttt{proof\_gaps/index.tsv} for
\texttt{1607.06127}, title keywords, \texttt{collapsed},
\texttt{almost uncollapsed}, \texttt{bounded}, and related terms; no matching
packet or attempt was found.

Web searches on June 14, 2026 used exact and close phrase searches for:
\begin{itemize}
\item collapsed, almost uncollapsed, bounded map, and Banach spaces;
\item the source title together with almost uncollapsed and collapsed;
\item the exact Problem 2.6 wording and shorter variants.
\end{itemize}
They found the source paper but no separate later paper explicitly answering
this problem. Novelty confidence is moderate rather than high because the
example is elementary and may be folklore.

\section*{Human review recommendation}

Send to human as a likely valid full answer to Problem 2.6. The main semantic
point to check is whether the author intended an unstated restriction to
cocycles or to another class of maps; under the printed statement, the circle
map answers the question.

\begin{thebibliography}{9}
\bibitem{Braga}
B. M. Braga,
\emph{On weaker notions of nonlinear embeddings between Banach spaces},
arXiv:1607.06127, 2016.
\end{thebibliography}

\end{document}
